\documentclass[11pt,a4paper]{article}
\usepackage[margin=25mm]{geometry}
\usepackage{amsmath,amssymb,amsthm,hyperref}
\setlength{\emergencystretch}{2em}
\newtheorem{theorem}{Theorem}
\newtheorem{corollary}[theorem]{Corollary}
\newcommand{\gen}[1]{\langle#1\rangle}
\title{Paired cuts for a selected set of colours}
\author{Human Compression project: L044 continuation}
\date{}
\begin{document}\maketitle
Let $H=\gen{s_0,\ldots,s_{r-1}}$ act transitively on a finite set
$U$ of size $m$. Each $s_i$ is an involution or the identity, and
$[s_i,s_j]=1$ for $|i-j|>1$. For every index $j$, put
$H_j=\gen{s_i:i\ne j}$. An edge of colour $j$ is eligible when
its endpoints lie in different orbits of this \emph{full local}
omission subgroup $H_j$.
Fix any colour set $F\subseteq\{0,\ldots,r-1\}$. It need not be
an interval, and $H_F=\gen{s_j:j\in F}$ need not be transitive.
Suppose every colour in $F$ has at least two eligible edges, and
select exactly two of each on the full vertex set $U$.

\begin{theorem}\label{main}
The bound $2|F|\leq m$ holds. If $2|F|=m$, there is a connected
paired selection using just the colours in $F$, and $H_F$ is
transitive on $U$. Consequently, if $H_F$ is intransitive, then
\[
                    2|F|\leq m-1.                 \tag{1}
\]
No intersection property, parity or primitivity is required.
\end{theorem}
\begin{proof}
We retain the full accepted proofs of \emph{Paired orbit cuts in
string actions} and \emph{Localisation in the paired-cut equality
case}. We check every place where selecting fewer than all colours
could change their argument.

Orbit-cut exchange, the selected-square block lemma and long-cycle
reduction use only the full local $H_j$ cuts and the fact that each
selected colour has exactly two edges. They therefore give a
square-only selection. Ordinary slides between squares use colours
of selected connector edges, so all colours they exchange belong
to $F$. If such a connector begins with the middle colour $i$ of
a square of colours $i-1,i+1$, then $i$ is itself selected. The
global middle-colour move retains this square and transports every
other selected edge by $s_i$. Every transported colour other than
$i$ commutes with $s_i$, except the two square colours, whose
selected edges have all been retained; $i$-edges map to themselves.
Eligibility is preserved by $s_i\in H_j$ for $j\ne i$.
Thus the complete minimum-square and minimum-distance proof still
gives at most one square per component.

If $c,q$ are the numbers of components and squares in such a
selection, including isolated vertices among the components, its
exact defect is $2|F|=m-c+q$. Hence $2|F|\leq m$. In equality,
every minimum-square selection has exactly one square in each
component. We now adapt the complete equality connectedness proof.

First, every eligible edge whose colour belongs to a selected square
lies in that square's component. An edge in another component
would have a connecting selected path avoiding its colour. An edge
joining components could replace a square edge, destroying that
square and joining components without a new cycle. These are the
two contradictions in the accepted cycle-colour localisation lemma.

Next consider a full graph edge of any colour $k$, selected or not,
incident with a square vertex. If $s_k$ commutes with a square colour
$j$, transporting its eligible $j$-edge by $s_k\in H_j$ and applying
localisation keeps the new endpoint in that component. The only
other case has square colours $i-1,i+1$ and $k=i$.
Here $i$ may be \emph{unselected}. Nevertheless the same global
exchange is valid: retain the square and transport all other selected
edges by $s_i$. If $i\notin F$ there are simply no selected
$i$-edges to account for. All other selected colours outside the
square commute with $s_i$, and their eligibility is preserved.

Viewed on the transported graph, the exchange removes the transported
square and inserts the original square. The removed square splits
its old unicyclic component into four trees. Every inserted edge
lies in the inserted full square, so creates no other simple cycle.
The number of squares is unchanged. Each inserted vertex must lie
in those four trees: if it lay in another unchanged component, the
new component would have two squares, contrary to the minimum-square
property. All four trees must be rejoined, or a tree component would
remain in equality. Thus the new component partition is exactly the
transported old partition, and the incident full edge stays in the
original square's component. This is precisely the earlier
full-neighbour and partition argument; the missing middle colour
causes no exception.

If the selection were disconnected, full $H$ transitivity would
give a cross-edge between components. Over all minimum-square
selections, choose a shortest connector consisting of selected tails
from two squares and exactly one full cross-edge. Both tails are
nonempty by the full-neighbour property. The accepted ordinary-slide
argument shortens either tail unless its first colour is the middle
colour of its square. Its repeated-colour exception is ruled out by
cycle-colour localisation, and ordinary slides preserve the component
vertex sets, exactly as before.

Both first colours must therefore be middle colours. The two squares
have disjoint colour sets. The cross-edge colour $x$ is outside the
colour set of at least one of them; choose that endpoint, with middle
colour $i$. Then $s_i$ commutes with $s_x$, including $x=i$.
The global middle exchange transports the cross-edge to an actual
$x$-edge and transports every connector edge after the first. The
transported component partition makes it a cross-edge still, and
the new tail starts on the retained square and is shorter by one.
The middle colour here belongs to $F$, since it was a selected tail
edge. This gives the same final contradiction as in the accepted
connectedness proof. Thus the equality selection is connected.
Its edges all have colours in $F$, so $H_F$ is transitive.
The integer strict bound (1) follows.
\end{proof}

\begin{corollary}[Exact local component interface]
Let $O$ be an orbit of a tail in an ambient string action, and let
$H$ be that tail's full restricted action on $O$. Suppose that two edges of every colour in $F$ cross
orbits of the actual full local omission $H_j$. Then
$|F|\leq |O|/2$. If some index $k\notin F$ has $H_k$
intransitive, then $|F|\leq(|O|-1)/2$.
\end{corollary}
\begin{proof}
Apply Theorem~\ref{main}, padding the other indices with identity
restrictions when needed. In the second case $H_F\leq H_k$ is
intransitive. No identification with an ambient omission is used.
\end{proof}

\paragraph{Attribution and limits.}
This supplies the selected-colour local rank interface appearing in
Cameron--Fernandes--Leemans--Mixer, \emph{Highest rank of a polytope
for $A_n$}, arXiv:1605.09173v1, Propositions 5.27 and 5.31,
physical pages 35--36 and 42--43. It does not prove that any particular
source component satisfies its local eligibility or intransitivity
hypotheses. Ambient omission eligibility and full local omission
eligibility are distinct: neither is silently substituted for the
other. It also asserts no sum of orbit deficits stronger than (1).
The root-supplied global middle move and its complete distinct
review remain bound through the paid predecessor manuscripts.
The present change is the explicit selected-colour scope argument.
There is no new census, formalization, source-wide recovery,
original credit or question credit. Connected-square rigidity and
the source's later component classification remain open obligations.
\end{document}
